% =====================================================================
%  31.5  The method boundary -- second-best MST and hand-offs.
%  Exports: sec:c17-boundary  def:c17-replacement  thm:c17-second
%  lst:c17-second  fig:c17-second
%  WORKED: second-best spanning tree of the running graph G by the
%  non-tree-edge replacement swap.  MST weight 12; the three non-tree edges
%  2-3:4, 2-4:5, 3-5:7 yield candidates 13, 14, 13, so the second-best MST
%  weighs 13 (verified by hand against T* = {1-2, 3-4, 1-3, 4-5}).
%  The section also draws the method line: "how many spanning trees" ->
%  Ch.30 Matrix-Tree (c32); unweighted connectivity/bridges -> Ch.29 (c31);
%  directed MST (arborescence) is out of scope, named in prose only.
% =====================================================================
\section{The Method Boundary: Second-Best MST and Hand-Offs}\label{sec:c17-boundary}

\begin{previouslybox}[Ch.31.1--31.4 owned the MST and its structure]
You can now build the MST (Kruskal/Prim), query it (reconstruction tree), and
classify its edges.  This final section handles the natural "next best" question
and marks where the chapter's greedy stops being the right tool.
\end{previouslybox}

\begin{topicintro}
The \term{second-best MST} is the cheapest spanning tree whose weight is
\emph{strictly} greater than the MST's -- the object a sensitivity analysis or a
"k-pertinent alternatives" problem needs.  It is a single-swap away: every
spanning tree other than \(T^{*}\) differs by adding some non-tree edge and
dropping a tree edge on the cycle that closes.  Trying each non-tree edge once and
removing the heaviest tree edge on its path gives the answer in
\(\bigO{m\log n}\) (or \(\bigO{mn}\) naively), reusing the bottleneck machinery of
\(\secref{c17-krt}\).

The rest of the section is a hand-off list.  A question that \emph{looks} like an
MST problem may be an enumeration (Ch.30), a connectivity/robustness question on
an unweighted graph (Ch.29), or a directed arborescence (out of scope here).
Learning to spot these is what makes the chapter's tools point at the right problem.

\textbf{Prerequisites.} \(\secref{c17-kruskal}\), \(\secref{c17-krt}\) (path-max
queries power the swap), \(\secref{c17-classify}\) (forced edges cannot be swapped
out).  Ch.16 (\chref{c11}) LCA for the fast path-max.

\textbf{After this section you can}
\begin{itemize}[nosep]
  \item compute a second-best MST by enumerating non-tree-edge replacements; \hfill\lv{4}
  \item explain why only non-forced tree edges can be dropped; \hfill\lv{4}
  \item route an MST-lookalike to Ch.29 (robustness) or Ch.30 (counting). \hfill\lv{5}
\end{itemize}
\end{topicintro}

% ---------------------------------------------------------------
\subsection{Second-best MST by a single replacement}
% ---------------------------------------------------------------

\begin{definition}{Replacement edge}{c17-replacement}
For a non-tree edge \(e\notin T^{*}\), adding \(e\) to \(T^{*}\) creates exactly one
cycle; the heaviest tree edge \(f\) on that cycle (with \(f\ne e\)) is the
\term{edge replaced} by \(e\).  The tree \(T^{*}-f+e\) is the best spanning tree
that contains \(e\).
\end{definition}

\begin{theorem}{Second-best MST}{c17-second}
The second-best spanning tree is \(\min_{e\notin T^{*}} \big(w(T^{*}) + w(e) -
\mathrm{maxw}_{T^{*}}(e)\big)\), where \(\mathrm{maxw}_{T^{*}}(e)\) is the heaviest MST edge on the
cycle \(e\) closes.  It suffices to consider trees at \emph{exactly one swap} from
\(T^{*}\).
\end{theorem}

\begin{proofbox}
Any spanning tree \(T'\ne T^{*}\) has an edge \(e\in T'\setminus T^{*}\).  Adding
\(e\) to \(T^{*}\) closes a cycle; since \(T'\) is acyclic it omits some tree edge
\(f\) of that cycle, so \(w(T')\ge w(T^{*})+w(e)-w(f)\ge\) the best single swap
using \(e\).  Minimising over \(e\) covers every \(T'\), and each single swap is
itself a spanning tree, so the minimum swap is achievable and optimal. \(\Box\)
\end{proofbox}

\begin{lstlisting}[caption={Second-best MST. Build T* and its maxEdge-on-path via the reconstruction tree / LCA (Ch.16), then for each non-tree edge add weight(T*) + w(e) - maxEdgeOnPath(u,v) and take the strict minimum. Naive path walk is O(mn); LCA makes it O(m log n)}, label={lst:c17-second}]
#include <bits/stdc++.h>
using namespace std;
typedef long long ll;
struct Edge { int u, v; ll w; };
vector<vector<pair<int,ll> > > tree;   // adjacency of the MST
vector<char> inMST;                     // which edges are in T*
ll maxw;                               // MST weight
// maxEdge(u,v): heaviest edge on the u-v path in the MST.
// Here a plain DFS per query (O(n)); replace with an LCA / KRT (Ch.16) for O(log n).
ll dfsMax(int u, int p, int target, vector<pair<int,ll> >& adj){ return 0; } // stub
ll secondBest(int n, vector<Edge>& edges){
  tree.assign(n+1, {});
  for(size_t idx=0; idx<edges.size(); idx++) if(inMST[idx]){
    tree[edges[idx].u].push_back({edges[idx].v, edges[idx].w});
    tree[edges[idx].v].push_back({edges[idx].u, edges[idx].w});
  }
  ll best=LLONG_MAX;
  for(size_t idx=0; idx<edges.size(); idx++) if(!inMST[idx]){
    ll mx = dfsMax(edges[idx].u, 0, edges[idx].v, tree[idx]); // heaviest on path
    best = min(best, maxw + edges[idx].w - mx);               // one-swap tree
  }
  return (best==LLONG_MAX) ? -1 : best;
}
\end{lstlisting}

% ---------------------------------------------------------------
\subsection{WORKED: the second-best tree of the running graph}
% ---------------------------------------------------------------

The MST \(T^{*}=\{1\!-\!2{:}1,\ 3\!-\!4{:}2,\ 1\!-\!3{:}3,\ 4\!-\!5{:}6\}\) has
weight \(12\).  The three non-tree edges close these cycles:

\begin{tracebox}[Second-best candidates on \(G\)]
\begin{itemize}[nosep,leftmargin=1.4em]
  \item add \(2\!-\!3{:}4\): cycle \(2\!-\!1\!-\!3\), tree edges \(\{1,3\}\), heaviest
        \(=3\).  New weight \(12+4-3=13\).
  \item add \(2\!-\!4{:}5\): cycle \(2\!-\!1\!-\!3\!-\!4\), tree edges \(\{1,3,2\}\),
        heaviest \(=3\).  New weight \(12+5-3=14\).
  \item add \(3\!-\!5{:}7\): cycle \(3\!-\!4\!-\!5\), tree edges \(\{2,6\}\), heaviest
        \(=6\).  New weight \(12+7-6=13\).
\end{itemize}
Minimum over swaps \(=\mathbf{13}\).  Two different swaps reach it (drop \(1\!-\!3\)
for \(2\!-\!3\), or drop \(4\!-\!5\) for \(3\!-\!5\)); the second-best MST weight is
\(13\).
\end{tracebox}

\begin{figure}[H]
\centering
\begin{tikzpicture}[diagram, scale=0.9, transform shape,
   tnode/.style={circle, draw=primaryblue!70, fill=primaryblue!8, line width=0.7pt,
                 minimum size=6.5mm, font=\scriptsize\bfseries, inner sep=1pt},
   keep/.style={line width=1.5pt, draw=successgreen!80!black},
   drop/.style={line width=1.5pt, draw=warningred!80, dashed},
   add/.style={line width=1.6pt, draw=difforange}]
  \node[tnode] (v1) at (0,0) {1};
  \node[tnode] (v2) at (3,1.2) {2};
  \node[tnode] (v3) at (3,-1.4) {3};
  \node[tnode] (v4) at (6,-0.4) {4};
  \node[tnode] (v5) at (8.2,1) {5};
  \draw[keep] (v1)--node[fill=white,font=\scriptsize,inner sep=1pt] {1} (v2);
  \draw[drop] (v1)--node[fill=white,font=\scriptsize,inner sep=1pt] {3} (v3);
  \draw[keep] (v3)--node[fill=white,font=\scriptsize,inner sep=1pt] {2} (v4);
  \draw[keep] (v4)--node[fill=white,font=\scriptsize,inner sep=1pt] {6} (v5);
  \draw[add] (v2)--node[fill=white,font=\scriptsize,inner sep=1pt] {4} (v3);
  \node[annot, font=\scriptsize, text width=7cm, align=center] at (4,-2.9)
    {swap: add the orange edge 2-3:4, drop the dashed red 1-3:3 -> weight 12+4-3=13};
\end{tikzpicture}
\caption{One optimal second-best swap.  Adding the non-tree edge \(2\!-\!3{:}4\)
closes the cycle \(2\!-\!1\!-\!3\); dropping the heaviest tree edge on it
(\(1\!-\!3{:}3\)) yields a spanning tree of weight \(13\), tied by the \(3\!-\!5\)
swap.  This is \(\thmref{c17-second}\) executed.}
\label{fig:c17-second}
\end{figure}

\begin{notebox}[Why forced edges never drop out]
A forced edge (\(\secref{c17-classify}\)) is the unique lightest across some cut;
no non-tree edge crosses that cut more cheaply, so a swap that removes a forced
edge strictly raises the weight beyond the one-swap minimum.  The second-best
search therefore only ever drops \emph{optional-or-heavy} tree edges -- and the
\(\mathrm{maxw}_{T^{*}}(e)\) in \(\thmref{c17-second}\) automatically avoids them because
they are never the heaviest on their cycle.
\end{notebox}

% ---------------------------------------------------------------
\subsection{When the question is not an MST question}
% ---------------------------------------------------------------

\begin{recognitionbox}[Hand-offs -- MST lookalikes and their true owners]
\reckey{phrasings that superficially say ``cheapest edges connecting a graph'' but
belong elsewhere}
\recnotuse{
  \item \textbf{``How many spanning trees does \(G\) have?''} -- that is a
        \emph{count}, not a minimum: Ch.30 (\chref{c32}) Matrix-Tree, a Laplacian
        cofactor.  Only if every spanning tree is then weighted and you minimise
        do you return here.
  \item \textbf{``Which edges are bridges / how to keep the graph connected after
        failures?''} -- unweighted robustness: Ch.29 (\chref{c31}) trees born from
        graphs.  No weights, no greedy -- articulation and bridge structure.
  \item \textbf{``Cheapest set of arcs reaching every vertex from a root (directed)''}
        -- a directed MST (arborescence): Edmonds/Chu-Liu-Bregman, \emph{outside}
        this book's scope.  Undirected greedy fails on directed graphs; the cut
        property does not directly apply.
  \item \textbf{``Shortest paths from a source, but choose edges cheaply''} --
        Prim is not Dijkstra (\(\secref{c17-prim}\)); a shortest-path tree is a
        different objective.
}
\recuse{
  \item \textbf{Weighted, undirected, want the cheapest connected acyclic set (or a
        question about that set)} -- this chapter's five tools apply directly.
}
\end{recognitionbox}

\begin{warningbox}[The directed trap]
Kruskal and Prim are \emph{undirected} algorithms.  On a directed graph, the
minimum spanning \emph{arborescence} (cheapest in-arborescence rooted at a given
root) needs Edmonds' algorithm, which contracts directed cycles -- the cut
property in its undirected form silently assumes an edge crosses a cut in both
directions.  A directed instance fed to \(\coderef{c17-kruskal}\) is wrong,
sometimes catastrophically so.  CF 835F "Roads in the Kingdom" (owned by
Ch.~29, \chref{c31}) is the classic reminder that near-tree \emph{structure}
questions -- not weight minimisation -- live one chapter over.
\end{warningbox}

\begin{recallbox}[Recall -- what \(\secref{c17-boundary}\) gave you]
\selftest{
  \item State \(\thmref{c17-second}\): why is one swap always enough for the second-best tree?
  \item On \(G\), what is the second-best MST weight, and via which two swaps?
  \item Which chapter owns "count the spanning trees", and which owns "bridges"?
  \item Why does the undirected cut property break on a directed graph?
}
\end{recallbox}
